\section*{Chapter 15}
\paragraph{15.A1.} Convolution distributes over addition and scalar
multiplication.
\paragraph{15.A2.} The output coordinate is always the same weighted pattern
of input coordinates, shifted cyclically.
\paragraph{15.B1.} Each row is a shifted copy of the kernel.
\paragraph{15.B2.} Multiply the circulant matrix by the vector $u$.
\paragraph{15.C1.} Use two nested loops and the index \texttt{(j - k) \% N}.
\paragraph{15.C2.} Compute the FFTs of $h$ and $u$, multiply them entrywise,
then apply the inverse FFT and take the real part.
\paragraph{15.D1.} Diagonal matrices act coordinate by coordinate; FFT makes
the needed basis change fast.
